\section{Introduction}
\label{sec:intro}

Coding agents built on large language models can now read a code base, run experiments, edit
proofs and write prose for days at a time. Whether they can carry a research program is a
different question, and the evidence so far comes mostly from short, benchmark-shaped tasks: an
idea to test, a paper to replicate, a Kaggle competition to enter
\citep{lu2024aiscientist,chan2025mlebench,starace2025paperbench}. Real research is longer, its
success criteria move, and its characteristic failure is not a crash but a claim that is slightly
stronger than what was shown. In numerical analysis this failure mode is acute: a convergence plot
can look sixth order for the wrong reason, a comparison with a public code can be unfair in ways
that take a domain expert to see, and a theorem can be true of a slightly different method than
the one implemented.

This paper is a longitudinal case study of one such program, conducted between May and September
2026 by LLM coding agents under the direction of a single human researcher. The subject of the
program was a high-order time integrator for constrained multibody systems modelled in absolute
coordinates on a Lie group. Its outcome is a method, \method{}, with an exact algebraic identity at
its core, a sixth-order error theorem whose algebra and perturbation lemmas are machine-checked in
Lean~4, and a numerical study on a frictional chain and on the four public ASME benchmark
mechanisms; that result is reported in a companion manuscript \citep{wang2026gauss6} and
summarized for a general reader in Figure~\ref{fig:problem}. Here we
report on the process: what the agent did, how the work was organized so that its claims could be
trusted, what went wrong, what caught it, and what the human actually contributed.

We make four contributions. (1)~\textbf{A documented instance} of agentic research with a
non-trivial result: the 49 versioned experiments, the 341 ledger scripts, the Lean development,
the manuscript and the transcript of the final phase are released, so the process can be inspected
rather than described. (2)~\textbf{A process architecture} (Section~\ref{sec:system}) of six
components that kept the program honest, from a versioned exploration tree to a manuscript
pipeline whose validator checks every printed number against its result file.
(3)~\textbf{A trajectory analysis} (Section~\ref{sec:trajectory}) of how the method emerged:
which hypotheses were rejected and why, and how an audit of the manuscript against the code
produced the central theorem. (4)~\textbf{An incident analysis} (Section~\ref{sec:incidents}) of
ten wrong or misleading claims, classified by what detected them: validators and purpose-built
experiments caught four, the agent's own cross-checks four, and the human two, both by looking at
a page or a figure and saying that it looked wrong.

Our reading of the evidence is that the agent was reliable at what can be checked mechanically and
unreliable at what cannot, and that the design consequence is not to keep the human in every loop
but to build loops that end in an artefact a machine can check, routing the few residual
judgments to the human.
